%!TEX program = xelatex
%!TEX encoding = UTF-8
\def\ano{2024}
\def\semestre{1}
\def\disciplina{Álgebra 1}
\def\turma{C}
\def\autor{José Antônio O. Freitas}
\def\instituto{MAT-UnB}

\documentclass{beamer}
\usetheme{Madrid}
\usecolortheme{beaver}
% \mode<presentation>
\usepackage[singlelinecheck=false]{caption}
\usepackage{lmodern}
\usepackage{textpos}
\usepackage{amssymb}
\usepackage{amsmath,amsfonts,amsthm,amstext}
\usepackage{cases}
\usepackage[justification=centering]{caption}
\usepackage[brazilian]{babel}
% \usepackage[latin1]{inputenc}
\usepackage{graphicx}
\graphicspath{{/home/jfreitas/GitHub/geometria_analitica/apresentacoes/logo}{/home/jfreitas/GitHub/geometria_analitica/notas_de_aula/figuras-latex}{C:/Users/jfreitas/GitHub/geometria_analitica/apresentacoes/logo}}
\usepackage{enumitem}
\usepackage{multicol}
\usepackage{answers}
\usepackage{qrcode}
\usepackage{tikz,ifthen}
\usetikzlibrary{lindenmayersystems}
\usetikzlibrary[shadings]

\newcounter{saveenumi}
\newcommand{\seti}{\setcounter{saveenumi}{\value{enumi}}}
\newcommand{\conti}{\setcounter{enumi}{\value{saveenumi}}}

\resetcounteronoverlays{saveenumi}


\renewcommand{\qedsymbol}{$\blacksquare$}

\setbeamertemplate{itemize items}[ball]

\input{../configuracoes_latex/ambientes}
\newcommand{\nsub}{\varsubsetneq}
\newcommand{\vaz}{\emptyset}
\newcommand{\dom}{{\rm dom\,}}
\newcommand{\im}{{\rm Im\,}}
\newcommand{\sub}{\subseteq}
\newcommand{\n}{\mathbb{N}}
\newcommand{\z}{\mathbb{Z}}
\newcommand{\rac}{\mathbb{Q}}
\newcommand{\real}{\mathbb{R}}
\newcommand{\complex}{\mathbb{C}}
\newcommand{\cp}[1]{\mathbb{#1}}
\newcommand{\ch}{\mbox{\textrm{car\,}}\nobreak}
\newcommand{\vesp}[1]{\vspace{ #1  cm}}
\newcommand{\compcent}[1]{\vcenter{\hbox{$#1\circ$}}}
\newcommand{\comp}{\mathbin{\mathchoice
{\compcent\scriptstyle}{\compcent\scriptstyle}
{\compcent\scriptscriptstyle}{\compcent\scriptscriptstyle}}}
\newcommand\T{\rule{0pt}{2.6ex}}
\newcommand{\elipse}[2]{\dfrac{x^2}{#1} + \dfrac{y^2}{#2} = 1}
\newcommand{\inner}[2]{\langle \vec{#1}, \vec{#2} \rangle}
\newcommand{\proj}{{\rm proj}}
\newcommand{\eq}{=}
\newcommand{\norm}[1]{\left\Vert{#1}\right\Vert}
\newcommand{\integer}{\mathbb{Z}}
\renewcommand{\sin}{{\rm sen\,}}
\renewcommand{\tan}{{\rm tg\,}}
\renewcommand{\csc}{{\rm cossec\,}}
\renewcommand{\cot}{{\rm cotg\,}}
\renewcommand{\sinh}{{\rm senh\,}}




\title{Retas no Plano}
\author[\autor]{\autor}
\institute[\instituto]{\instituto}
\date{}

\begin{document}
\begin{frame}
    \maketitle
\end{frame}

\logo{
\includegraphics[scale=.1]{logo-MAT.png}\vspace*{8.5cm}}

\begin{frame}
    \begin{definicao}
        Qualquer vetor não nulo \pause e paralelo \pause a uma dada reta $r$ \pause é chamado de \textbf{vetor diretor} de $r$.
    \end{definicao}
\end{frame}

\begin{frame}
    Seja $\vec{u}$ um vetor diretor de uma reta $r$ \pause e seja $A$ um ponto pertencente à reta $r$. \pause Assim um ponto $X$ \pause pertence a reta $r$
    se, \pause e somente se, \pause $\vec{AX}\varparallel\vec{u}$.\pause
    \begin{figure}[!h]
        \centering
        \caption{Equação vetorial da reta}
        \includegraphics{equacao-vetorial-reta-plano-cartesiano.pdf}
    \end{figure}
\end{frame}

\begin{frame}
    \begin{definicao}\index{Reta no Plano!Equação Vetorial}
        A equação\pause
        \[
            \vec{AX} = \lambda \vec{u},\pause
        \]
        onde $\lambda \in \real$ \pause é chamada de \textbf{equação vetorial \pause da reta} $r$ \pause ou \textbf{equação da reta} $r$ \pause
        \textbf{na forma vetorial}.
    \end{definicao}
\end{frame}

\begin{frame}
    \begin{observacao}
        Como o ponto \pause e o vetor diretor \pause de uma reta $r$ \pause são escolhidos de forma arbitrária, \pause existem infinitas equações vetoriais para uma
        mesma reta $r$. \pause Por exemplo, \pause se $A$ e $B$ são pontos distintos de uma reta $r$, \pause então $\vec{AB}$ \pause e $\vec{BA}$ \pause são vetores diretores distintos de $r$ e portanto temos as equações vetoriais\pause
        \begin{align*}
            \vec{AX} = \lambda\vec{AB} \\
            \vec{AX} = \alpha\vec{BA}.
        \end{align*}
        Na verdade, \pause qualquer vetor múltiplo do vetor $\vec{AB}$ \pause pode ser tomado como vetor diretor de $r$.
    \end{observacao}
\end{frame}

\begin{frame}
    \begin{definicao}
        Seja $r$ uma reta \pause com vetor diretor $\vec{u} = (a,b)$, \pause e que passa pelo ponto $(x_0, y_0)$. \pause O sistema de equações \pause
        \[
            \begin{cases}
                x = x_0 + \lambda a \\
                y = y_0 + \lambda b
            \end{cases}, \lambda \in \real
        \]\pause

        é chamado \textbf{sistema de equações paramétricas da reta} $r$ \pause ou \textbf{sistema de equações da reta} \pause $r$ \textbf{na forma paramétrica}.
    \end{definicao}
\end{frame}

\begin{frame}
    \begin{definicao}
        Seja $r$ uma reta \pause com vetor diretor $\vec{u} = (a,b)$, \pause onde $ab\ne 0$ \pause e que passa pelo ponto $(x_0, y_0)$. \pause A equação\pause
        \begin{equation*}
            \dfrac{x - x_0}{a} = \dfrac{y - y_0}{b}\pause
        \end{equation*}
        é chamada de \pause \textbf{sistema de equações da reta} $r$ \pause \textbf{na forma simétrica} \pause ou por abuso de linguagem, \pause
        \textbf{equações da reta} $r$ \pause\textbf{na forma simétrica}.
    \end{definicao}
\end{frame}

\begin{frame}
    \begin{definicao}
        Seja $r$ uma reta que contém o ponto $(x_0, y_0)$ \pause e com vetor diretor $\vec{u} = (a, b)$. \pause A equação \pause
        \begin{equation*}
            ay - bx = c\pause
        \end{equation*}
        onde $c = ay_0 - bx_0$ \pause é chamada de \pause \textbf{equação cartesiana da reta} $r$.
    \end{definicao}
\end{frame}

\begin{frame}
    \begin{figure}[!h]
        \centering
        \caption{ângulo entre retas}
        \includegraphics{angulo-retas-plano-cartesiano.pdf}
    \end{figure}
\end{frame}

\begin{frame}
    \begin{definicao}
        A distância de um ponto $P(x_0,y_0)$ \pause à reta $r$ de equação $y = mx + k$ \pause é definida como sendo a distância de $P$ \pause ao ponto $A(x_1,
            y_1)$, \pause onde $A$ é o pé da perpendicular \pause à reta $r$ \pause passando pelo ponto $P$.
    \end{definicao}
\end{frame}

\begin{frame}
    \begin{figure}[!h]
        \centering
        \caption{Distância de ponto a reta}
        \includegraphics{distancia-ponto-reta-plano-cartesiano.pdf}
    \end{figure}
\end{frame}

\end{document}
